Via NetworkManager kunnen we specifieke ethernet settings zetten op een interface.

We beginnen met het maken van een nieuwe interface. We geven de interface een eigen naam:
\begin{lstlisting}[language=bash]
$ sudo nmcli con add con-name "localethernet" type 802-3-ethernet
\end{lstlisting}

Van de vele settings die er mogelijk zijn geven we een paar voorbeelden.

Normaal gesproken wordt er op een ethernet link onderhandeld tussen de switch en de machine welke snelheden er gebruikt kunnen worden en of er half- of full-duplex mogelijk is. De parameter die hiervoor gebruikt wordt is \texttt{802-3-ethernet.auto-negotiate}. We kunnen deze op \textbf{false} zetten als we zelf willen bepalen hoe de verbinding eruit komt te zien.
\begin{lstlisting}[language=bash]
$ sudo nmcli con modify "localethernet" 802-3-ethernet.auto-negotiate <BOOLEAN>
\end{lstlisting}

Met de \texttt{802-3-ethernet.speed} parameter kunnen we de snelheid van de interface zetten in Mbits/s. Standaard is deze waarde unset (0), zodat hij onderhandeld (auto-negotiate) kan worden. Zetten we deze op een bepaalde waarde dan zal die waarde meegegeven worden in de negotiate, als auto-negotiate op \textbf{true} staat. Is auto-negotiate \textbf{false} is de snelheid de opgegeven snelheid.
\begin{lstlisting}[language=bash]
$ sudo nmcli con modify "localethernet" 802-3-ethernet.speed <INT>
\end{lstlisting}

Half- of full-duplex in ethernet betekent dat er twee of vier-draden tegelijk gebruikt kunnen worden. Half-duplex is dat er of gezonden of ontvangen kan worden, dus of de transmit(TX) of de receive(RX) draden worden gebruikt. Full-duplex betekent dat er tegelijktijd gezonden en ontvangen kan worden, de transmit(TX) en de receive(RX) draden worden tegelijkertijd gebruikt. Is deze setting \textbf{unset}, dan bepaald auto-negotiate de uitkomst, is er \textbf{half} of \textbf{full} gezet dan wordt deze waarde aangenomen.
\begin{lstlisting}[language=bash]
$ sudo nmcli con modify "localethernet" 802-3-ethernet.duplex <STRING>
\end{lstlisting}
Om de setting op \textbf{unset} te zetten verwijder hem uit het profiel.

